% GENERATED FILE. Edit canonical lessons, then run npm run generate:books.
% canonical-source-hash: fnv1a64:5ff6ced4296e2245
% canonical-lessons: HI-C276-vahgaya, HI-C276-usnekhai, HI-R276-what-happened

\chapter{What Happened, and Who Did It}
\label{ch:hi-what-happened-and-who-did-it}
\hlchaptermodality{276}

\begin{chapteropening}
\textbf{By the end of this chapter:} \emph{I can say what happened (vah gayā, vah āī) and who did something, with ne and a verb that follows the thing (usne roṭī khāī).}

Complete HI-R276-what-happened: say what happened and who did it, in the past.
\end{chapteropening}

\section[\texorpdfstring{\hi{वह} \hi{गया} (vah gayā)}{वह गया (vah gayā)}]{\hi{वह} \hi{गया} (vah gayā) — the past, by who did it}

\label{lesson:HI-C276-vahgaya}



\begin{quote}
Before the new one: say turning over and a stretch. (\emph{karvaṭ}, \emph{aṅgṛāī}.) Then the host's question \textbf{\hi{खाना} \hi{खाया}?}, \textquotedblleft{}have you eaten?\textquotedblright{} Its \textbf{\hi{खाया}} is a past, and this chapter builds every past like it.
\end{quote}

\subsection*{You'll want to know: the past, by who did it}
\begin{quote}
\textbf{\hi{कल} \hi{मैं} \hi{बाज़ार} \hi{गया।}}
\end{quote}

\begin{quote}
\emph{kal maiṁ bāzār gayā.} — \textquotedblleft{}Yesterday I went to the market.\textquotedblright{}
\end{quote}

With a past, \textbf{\hi{कल}} is yesterday. \textbf{\hi{गया}} is \textbf{\hi{जाना}} in the past, with no \textbf{\hi{है}} after it. Its ending follows the one who went, as \textbf{\hi{अच्छा}}, \textbf{\hi{अच्छी}}, \textbf{\hi{अच्छे}} follow a noun:

\noindent
\begin{tabularx}{\linewidth}{@{}>{\raggedright\arraybackslash}X>{\raggedright\arraybackslash}X@{}}
\toprule
\textbf{who went} & \textbf{the past} \\
\midrule
a man: \textbf{\hi{मैं}}, \textbf{\hi{वह}} & \textbf{\hi{गया}} (\emph{gayā}) \\
a woman & \textbf{\hi{गई}} (\emph{gaī}) \\
men, or \textbf{\hi{आप}}, \textbf{\hi{तुम}}, \textbf{\hi{हम}} & \textbf{\hi{गए}} (\emph{gae}) \\
women & \textbf{\hi{गईं}} (\emph{gaīṁ}) \\
\bottomrule
\end{tabularx}

\subsection*{You'll want to know: how a stem takes the ending}
Most verbs put the ending on the stem: \textbf{\hi{बैठना}} $\to$ \textbf{\hi{वह} \hi{बैठा}}, \textbf{\hi{वह} \hi{बैठी}} (\emph{vah baiṭhā}, \emph{vah baiṭhī}). A stem ending in a vowel adds \textbf{-\hi{या}} for a man: \textbf{\hi{आना}} $\to$ \textbf{\hi{आया}}, \textbf{\hi{आई}} (\emph{āyā}, \emph{āī}). Two change the stem: \textbf{\hi{जाना}} $\to$ \textbf{\hi{गया}}, and \textbf{\hi{होना}} $\to$ \textbf{\hi{हुआ}}, \textbf{\hi{हुई}} (\emph{huā}, \emph{huī}): \textbf{\hi{क्या} \hi{हुआ}?}, \textquotedblleft{}what happened?\textquotedblright{}

\textquotedblleft{}Did not\textquotedblright{}: \textbf{\hi{मैं} \hi{नहीं} \hi{गया}} (\emph{maiṁ nahīṃ gayā}). A bare stem with \textbf{\hi{गया}} says the change is complete, as in \textbf{\hi{देर} \hi{हो} \hi{गई}} and \textbf{\hi{वह} \hi{सो} \hi{गया}} (\emph{vah so gayā}), \textquotedblleft{}he fell asleep\textquotedblright{}.

\subsection*{Your turn}
Say these aloud:

\begin{itemize}
\raggedright
  \item \textquotedblleft{}She came home in the evening\textquotedblright{} — \emph{vah śām ko ghar āī}
  \item \textquotedblleft{}We sat down. Then we slept.\textquotedblright{} — \emph{ham baiṭhe. phir ham soe.}
\end{itemize}

\begin{tcolorbox}[breakable,colback=teal!4,colframe=teal!35!black,title={Before you move on}]
\textquotedblleft{}He came\textquotedblright{}? (\textbf{\hi{आया}}.) \textquotedblleft{}She went\textquotedblright{}? (\textbf{\hi{गई}}.) \textquotedblleft{}What happened?\textquotedblright{} (\textbf{\hi{क्या} \hi{हुआ}?}) \textquotedblleft{}I didn't go\textquotedblright{}? (\textbf{\hi{मैं} \hi{नहीं} \hi{गया}}.)
\end{tcolorbox}
\section[\texorpdfstring{\hi{उसने} \hi{रोटी} \hi{खाई} (usne roṭī khāī)}{उसने रोटी खाई (usne roṭī khāī)}]{\hi{उसने} \hi{रोटी} \hi{खाई} (usne roṭī khāī) — the past that follows the thing}

\label{lesson:HI-C276-usnekhai}



\begin{quote}
Say \textquotedblleft{}she went\textquotedblright{}. (\emph{vah gaī}.) \textbf{\hi{खाना} \hi{खाया}?} leaves the eater out; this lesson puts the eater in.
\end{quote}

\begin{grammarlens}[title={Grammar Lens: the doer takes \hi{ने}}]
\begin{quote}
\textbf{\hi{उसने} \hi{खाना} \hi{खाया।}}
\end{quote}

\begin{quote}
\emph{usne khānā khāyā.} — \textquotedblleft{}He (or she) ate the meal.\textquotedblright{}
\end{quote}

With a verb that has an object, such as eating, reading or saying, the past changes: the doer takes \textbf{\hi{ने}} (\emph{ne}). The pronouns fuse with it: \textbf{\hi{मैंने}}, \textbf{\hi{तुमने}}, \textbf{\hi{आपने}}, \textbf{\hi{उसने}}, \textbf{\hi{हमने}}, \textbf{\hi{उन्होंने}} (\emph{maiṁne}, \emph{tumne}, \emph{āpne}, \emph{usne}, \emph{hamne}, \emph{unhoṁne}), the last from \textbf{\hi{वे}}. A name takes \textbf{\hi{ने}} after it: \textbf{\hi{मीरा} \hi{ने}}. \textbf{\hi{उसने}} is \textquotedblleft{}he\textquotedblright{} or \textquotedblleft{}she\textquotedblright{}; the verb no longer says which.
\end{grammarlens}

\subsection*{You'll want to know: the verb follows the thing}
\begin{quote}
\textbf{\hi{मीरा} \hi{ने} \hi{रोटी} \hi{खाई।}}
\end{quote}

\begin{quote}
\emph{Mīrā ne roṭī khāī.} — \textquotedblleft{}Mira ate a roti.\textquotedblright{}
\end{quote}

\textbf{\hi{रोटी}} is feminine, so \textbf{\hi{खाई}}, whoever ate it. With \textbf{\hi{खाना}}, masculine, \textbf{\hi{खाया}}; with books, \textbf{\hi{मैंने} \hi{किताबें} \hi{पढ़ीं}} (\emph{maiṁne kitābeṁ paṛhīṁ}). When the thing takes \textbf{\hi{को}}, or there is none, the verb stays \textbf{-\hi{आ}}: \textbf{\hi{उसने} \hi{मीरा} \hi{को} \hi{देखा}} (\emph{usne Mīrā ko dekhā}), \textquotedblleft{}he saw Mira\textquotedblright{}.

\noindent
\begin{tabularx}{\linewidth}{@{}>{\raggedright\arraybackslash}X>{\raggedright\arraybackslash}X@{}}
\toprule
\textbf{sentence} & \textbf{the verb follows} \\
\midrule
\textbf{\hi{वह} \hi{गई}} — no thing & the one who went \\
\textbf{\hi{उसने} \hi{रोटी} \hi{खाई}} — \textbf{\hi{ने}} and a thing & the thing \\
\textbf{\hi{उसने} \hi{मीरा} \hi{को} \hi{देखा}} — the thing has \textbf{\hi{को}} & nothing: \textbf{-\hi{आ}} \\
\bottomrule
\end{tabularx}

Four verbs make this past their own way, for a masculine thing and a feminine one: \textbf{\hi{किया}} / \textbf{\hi{की}}, \textbf{\hi{दिया}} / \textbf{\hi{दी}}, \textbf{\hi{लिया}} / \textbf{\hi{ली}}, \textbf{\hi{पिया}} / \textbf{\hi{पी}} (\emph{kiyā} / \emph{kī}, \emph{diyā} / \emph{dī}, \emph{liyā} / \emph{lī}, \emph{piyā} / \emph{pī}), from \textbf{\hi{करना}}, \textbf{\hi{देना}}, \textbf{\hi{लेना}} and \textbf{\hi{पीना}}.

\textbf{\hi{कहना}} (\emph{kahnā}), \textquotedblleft{}to say\textquotedblright{}, is regular. Your \textbf{\hi{क्या} \hi{कहा}?} was its past: \textbf{\hi{उसने} \hi{कहा} \hi{कि} \hi{वह} \hi{कल} \hi{आएगा}} (\emph{usne kahā ki vah kal āegā}), \textquotedblleft{}he said he would come tomorrow\textquotedblright{}. \textbf{\hi{लाना}}, \textquotedblleft{}to bring\textquotedblright{}, does not take \textbf{\hi{ने}}: \textbf{\hi{वह} \hi{फल} \hi{लाया}} (\emph{vah phal lāyā}).

\subsection*{Your turn}
Say these aloud:

\begin{itemize}
\raggedright
  \item \textquotedblleft{}Mira drank tea\textquotedblright{} — \emph{Mīrā ne cāy pī}
  \item \textquotedblleft{}She saw Arun\textquotedblright{} — \emph{usne Aruṇ ko dekhā}
\end{itemize}

\begin{tcolorbox}[breakable,colback=teal!4,colframe=teal!35!black,title={Before you move on}]
Why \textbf{\hi{खाई}} in \textbf{\hi{उसने} \hi{रोटी} \hi{खाई}}? (\textbf{\hi{रोटी}} is feminine.) Now the other way: \textquotedblleft{}he ate the meal\textquotedblright{}? (\textbf{\hi{उसने} \hi{खाना} \hi{खाया}}.) Why \textbf{\hi{देखा}} in \textbf{\hi{उसने} \hi{मीरा} \hi{को} \hi{देखा}}? (The thing has \textbf{\hi{को}}.) \textquotedblleft{}He said\textquotedblright{}? (\textbf{\hi{उसने} \hi{कहा}}.)
\end{tcolorbox}
\section[(review)]{Review — what happened, and who did it}

\label{lesson:HI-R276-what-happened}



\begin{quote}
Say \textquotedblleft{}Yesterday I went to the market\textquotedblright{}. (\emph{kal maiṁ bāzār gayā}.) Then \textquotedblleft{}she came\textquotedblright{} and \textquotedblleft{}what happened?\textquotedblright{} (\emph{vah āī}, \emph{kyā huā?})
\end{quote}

\subsection*{Your turn}
Say these aloud:

\begin{itemize}
\raggedright
  \item I came home in the evening, and then I slept.
  \item Mira ate a roti. Arun ate the meal.
  \item What did you do yesterday? I went to the market.
  \item She saw Arun at the station.
  \item He said that he would come tomorrow.
\end{itemize}

\begin{tcolorbox}[breakable,colback=teal!4,colframe=teal!35!black,title={Before you move on}]
Which past follows the one who went? (A verb with no thing: \textbf{\hi{वह} \hi{गई}}.) Which follows the thing? (The past with \textbf{\hi{ने}}: \textbf{\hi{मीरा} \hi{ने} \hi{रोटी} \hi{खाई}}.) Now the other way round: in \textbf{\hi{उसने} \hi{खाना} \hi{खाया}}, which word tells you that \textbf{\hi{खाना}} is masculine? (\textbf{\hi{खाया}}.) And when the thing takes \textbf{\hi{को}}? (The verb stays \textbf{-\hi{आ}}: \textbf{\hi{उसने} \hi{मीरा} \hi{को} \hi{देखा}}.)
\end{tcolorbox}
